\prefacesection{Acknowledgments}

A dissertation is, in the end, a record of years of accumulated debts. The work presented here was made possible by a remarkable community of mentors, collaborators, lab-mates, staff, friends, and family, and it is a privilege to acknowledge some of them by name.

My deepest thanks go to my advisor, Anthony Wagner. Anthony took a chance on me when I arrived at Stanford with more enthusiasm than focus, and he has spent the years since teaching me, by patient example, what it looks like to do careful, generous, and ambitious science. His taste for the right problem, his insistence on precision in both thought and prose, and his characteristic refusal to accept a sloppy claim --- however well-marketed --- have shaped this dissertation more than any single paper or technique I learned along the way. Beyond the science, Anthony made room for the human side of training: the conversations about life and family, the encouragement during setbacks, the consistent reminder that the work matters but the people doing it matter more. I am profoundly grateful for his mentorship and his trust, and I will spend the rest of my career trying to pay forward what he has gifted me.

I am grateful to the members of my dissertation committee, Beth Mormino, Judy Fan, Kalanit Grill-Spector, and Russ Poldrack, whose careful reading and thoughtful questions sharpened the arguments in this document considerably. Each of them brought a distinct perspective to the table, and the dissertation is the better for it.

The empirical and methodological chapters of this dissertation are first-author papers, but they are not solo work. Haopei Yang (HY), Alice Xue, and Mingjian (Alex) He were the kind of collaborators every graduate student hopes for: rigorous, generous with their time, willing to argue, and unfailingly kind in the process. HY has been a co-conspirator across every project in this dissertation --- co-first-authoring the \textit{Attending to Remember} review, helping shape the \texttt{eyeris} preprocessing pipeline, and offering thoughtful counsel at nearly every stage of the PRISM work --- and I cannot imagine my time in the lab without him. Alice brought to the \texttt{eyeris} and attention--memory work the kind of clarity of thought and care for detail that make collaborators better just by working alongside them. Alex's deep expertise in signal processing and pupillometry methods elevated the \texttt{eyeris} package well beyond what I could have built alone. Alex, thank you for all of the late-night Slack conversations; your curiosity dramatically shaped my graduate training and the work that resulted from it. I am also grateful to Douglas Miller, Tammy Tran, Jintao Sheng, Natalie Biderman, and Subbulakshmi S for conceptual input, Khanh Nguyen for data collection, and to Kevin Madore and Megan deBettencourt for foundational conversations about closed-loop attention work that informed Chapters~\ref{ch:biofeedback-framework} and~\ref{ch:prism}.

The Stanford Memory Lab has been my intellectual home for the duration of this PhD, and the community that filled it day to day was, more than any single experiment, what made the training worth it. To the current lab members --- Douglas, Alice, Alex, HY, Subbu, Jintao, and Ali  --- and to those who came before me and went on to their next chapters --- Tyler Bonnen and Marc Harrison --- thank you for the journal clubs, the lab meetings, the hallway conversations that turned into proposals, the proposals that turned into papers, and the steady reminder that science is a team sport. The friendships made here are, I suspect, the most enduring thing I will take from graduate school.

Beyond my immediate lab, several faculty and postdocs at Stanford and adjacent institutions shaped my thinking in ways large and small. I am especially grateful to Laura Carstensen, Carol Dweck, Cameron Ellis, Mike Frank, Tobi Gerstenberg, Brian Knutson, Jay McClelland, and Tony Norcia for formative courses, informal advising, and collaboration. I also want to thank the FriSem community participants, whose questions at presentations and informal discussions pushed me to think more carefully about the relationship between measurement and inference that runs through this dissertation.

A separate and heartfelt thanks goes to Professor Judy Fan, who served as an unofficial secondary advisor over the course of my graduate training but who became, much more importantly, a close friend. I had the good fortune of being among the first people Judy spoke with when she arrived at Stanford as new faculty, and what began with those early conversations grew into a sustained collaboration on the cognitive science of learning and attention in real undergraduate classrooms, as well as into the larger project of building Psych 10 (Intro Stats) into the course it has become today. The hours spent in her office, in lecture halls, and over coffee --- debating how to teach statistics well and what it really means to support students through quantitative training --- are among the most defining experiences of my time in graduate school. The work we built together, and the friendship that grew alongside it, formed memories I will carry for the rest of my life.

It has been one of the genuine privileges of my time at Stanford to work alongside the undergraduate research assistants and honors-thesis students who trusted me to help guide their first steps into research. Thank you for the questions I would not have known to ask, for your patience as I learned how to mentor in real time, and for the reminder, on the days when my own work felt stuck, that the most rewarding part of training is what gets passed forward. I have learned at least as much from working with you as I have managed to teach.

To the several hundred students who passed through Psych 10 and Psych 45 (Learning \& Memory) while I served as their head teaching assistant --- thank you. Standing in front of a section full of curious, skeptical, and thoughtful undergraduates, week after week, shaped my own thinking about how to communicate science (and what is worth communicating in the first place). You were generous with your engagement, gracious with my early-career missteps, and a recurring reminder that the eventual audience for this work is people who want to understand the world a little better than they did the week before.

I am also indebted to the broader Stanford Psychology Department community --- faculty, postdocs, graduate students, and undergraduates alike --- for trusting me to run the department's Stats Consulting service and Hacky Hours for the past several years. Some of the most engaging and intellectually stretching hours of my time at Stanford were spent across the table from colleagues wrestling with the statistical questions their data had thrown at them. Every question that walked through the door was, in some sense, an invitation to grow, and I always looked forward to them. Whatever statistical fluency I leave Stanford with was built, in no small part, by trying to be useful to you. A huge thank you to Julia Proshan and Catherine Garton for helping make Stats Hacky Hours a mainstay in the department.

I would also like to briefly thank the various committees in the Psychology Department and at the Wu Tsai Neurosciences Institute on which I had the chance to serve during graduate school. I am grateful to have been trusted with the opportunity to shape, in small ways, the institutional environment that current and future trainees will inherit.

Long before Stanford, a series of patient and generous mentors made graduate school possible in the first place. I owe particular thanks to Alan Castel, who introduced me to research and to the discipline of asking questions with real care, and to Michael Alfaro, who prepared me to apply to and ultimately complete my PhD. I would not have known how to begin without them.

Science depends on the largely invisible labor of the staff who keep our programs and our buildings running. I am especially grateful to Dena Zlatunich, Justin Dixon, Andrea Sims, Diana Savastio, and Harry Bahlman for the administrative work they do to keep the lights on (figuratively \& literally). Additionally, I am grateful to the Stanford Memory Lab Managers, past and present --- Hannah Lloyd, Julia Rathmann-Bloch, America Romero, Jen Park, Gloria Cheng, Khanh Nguyen, Lucah Medina Guerra, Isha Sai, and Tyler Ward --- who kept participant recruitment and lab logistics running smoothly over the years. Without their steady, often thankless work, none of the studies reported here would have happened.

To my graduate cohort and the broader community of friends I made during my time at Stanford --- thank you for the friendship that outlasted any given academic quarter. Particular thanks to Josh Wilson, Sarah Tung, and Julia Proshan, who kept me grounded when the work got heavy and celebrated with me when it came together. To friends from earlier seasons of my life --- Hussain Ali and Hannah Nizinski --- thank you for staying close across the distance and the years.

To my wife, Alma, who has lived this dissertation as closely as I have and who stood beside me through some of the most difficult chapters of my life --- not only the PhD itself but the personal and health-related seasons that ran alongside it. This dissertation is as much yours as it is mine. Through every long stretch when the work demanded everything I had, and through every harder turn outside of the lab, you held the line. You picked up the slack at home when I was unavailable, absorbed the late nights and the missed dinners and the weekends I disappeared into a writing window, and never once let me believe that I was facing any of it alone. Your patience and understanding were not a passive gift; they were a daily, deliberate choice, made again and again --- without hesitation, without resentment, and without ever walking away. That you made that choice through the hardest moments of these years, when no one would have begrudged you the alternative, is the single greatest grace of my adult life.

Above all, you are my best friend. You are the person I most want to tell when something good happens and the person whose perspective I most trust when something hard does. The richness of these years --- the parts that will outlast this document --- comes from the life we have built together: the small things, the inside jokes, the routines, the quiet weeknights, the restful weekends, the way you make a house feel like home. Whatever else I take from this PhD, by far the most important is you, and the life we get to keep building. I love you.

And to my emotional support corgi, Reese, whose patient supervision of every long writing day and joyful insistence on regular outdoor breaks did more for the completion of this dissertation than he will ever know. Some of the clearest thinking I did during these years was on a walk, with him. I could not have gotten through this PhD without him. I love you. 

\bigskip

\noindent This work was supported, in part, by a Stanford University Wu Tsai Human Performance Alliance Agility Grant (to S. T. Schwartz \& A. D. Wagner); the Stanford University Wu Tsai Neurosciences Institute Center for Mind, Brain, Computation and Technology (to S. T. Schwartz \& A. M. Xue); a Stanford University Ric Weiland Graduate Fellowship in the Humanities and Sciences (to S. T. Schwartz); the McKnight Brain Research Foundation Clinical Translational Research Scholarship in Cognitive Aging and Age-Related Memory Loss; the American Brain Foundation; the American Academy of Neurology (to H. Yang); National Institutes of Health Grant R01-AG065255 (to A. D. Wagner); and National Science Foundation Graduate Research Fellowship DGE-2146755 (to A. M. Xue).

\bigskip

\noindent Chapter~\ref{ch:attending-to-remember} is a reprint of: Schwartz, S. T., Yang, H., Xue, A. M., \& Wagner, A. D. (2025). Attending to remember: Recent advances in methods and theory. \textit{Current Directions in Psychological Science}, 34(6), 330--341. \href{https://doi.org/10.1177/09637214251339452}{doi:10.1177/09637214251339452}. S. T. Schwartz and H. Yang contributed equally to this work. All authors approved the final manuscript for submission.

\noindent Chapter~\ref{ch:eyeris} is a reprint of: Schwartz, S. T., Yang, H., Xue, A. M., \& He, M. (2025). \texttt{eyeris}: A flexible, extensible, and reproducible pupillometry preprocessing framework in {R}. \textit{bioRxiv}. \href{https://doi.org/10.1101/2025.06.01.657312}{doi:10.1101/2025.06.01.657312}. All authors approved the final manuscript for submission.

\noindent Chapter~\ref{ch:biofeedback-framework} is reproduced from an unpublished manuscript drafted by S. T. Schwartz during the early years of doctoral training. It has not been previously submitted for publication.

\noindent Chapter~\ref{ch:prism} is in preparation for submission as a co-authored manuscript: Schwartz, S. T., Yang, H., He, M., Nguyen, K. K., deBettencourt, M. T., Tran, T. T., Madore, K. P., \& Wagner, A. D. (in preparation). \textit{Real-time reorienting during episodic retrieval rescues memories from suboptimal arousal-attentional states}. Author contributions: S.T.S.: Conceptualization, Formal Analysis, Funding Acquisition, Methodology, Investigation, Project Administration, Software, Writing --- Original Draft Preparation; H.Y.: Conceptualization, Formal Analysis, Methodology, Supervision, Writing --- Review \& Editing; M.H.: Formal Analysis, Supervision, Writing --- Review \& Editing; K.K.N.: Investigation, Project Administration, Writing --- Review \& Editing; M.T.dB.: Methodology, Writing --- Review \& Editing; T.T.T.: Supervision, Writing --- Review \& Editing; K.P.M.: Conceptualization, Funding Acquisition, Writing --- Review \& Editing; A.D.W.: Conceptualization, Formal Analysis, Funding Acquisition, Methodology, Supervision, Writing --- Review \& Editing.
